% A proved linear-algebra obstruction is written here for independent audit.
% No generic alternating-form rank or positive-index rank is asserted.
\subsection{Failure of critical Frobenius-fibre interpolation on products}

\begin{lemma}[Two Frobenius fibres impose seven conditions]
\label{lem:DE-A4-product-Frobenius-interpolation}
\status{Proved; complete independent root and monodromy audits passed.}
Let $k$ be algebraically closed of characteristic two. Let $E_1,E_2$
be elliptic curves, let $L_i$ have degree three on $E_i$, and set
\[
X=E_1\times E_2,\qquad L=L_1\boxtimes L_2,\qquad
\mathcal E=F_{X/k*}L.
\]
Then $h^0(X,L)=9$ and $\mathcal E$ is a rank-four bundle on $X^{(2)}$.
For two points $q=(q_1,q_2)$ and $q'=(q'_1,q'_2)$ of $X^{(2)}$
with $q_i\ne q'_i$ for both $i$, the evaluation map
\[
H^0(X^{(2)},\mathcal E)\longrightarrow
                         \mathcal E_q\oplus\mathcal E_{q'}
\tag{DE-FI1}
\]
has rank exactly seven. If the two points are distinct but one
coordinate agrees, its rank is six. In particular, evaluation at
two general points is not surjective, although the target has
dimension eight and $2=(h^0(L)-1)/4$.
\end{lemma}
\begin{proof}
For a point $q_i\in E_i^{(2)}$, write $p_i$ for its unique geometric
preimage. The Frobenius fibre is the length-two divisor $2p_i$.
Finite flat base change identifies the fibre of
$\mathcal E_i=F_{E_i/k*}L_i$ with
$H^0(2p_i,L_i|_{2p_i})$. The exact restriction sequence is
\[
0\longrightarrow L_i(-2p_i)\longrightarrow L_i
          \longrightarrow L_i|_{2p_i}\longrightarrow0.
\]
On an elliptic curve a degree-one line bundle has one section and
no first cohomology: its dual has negative degree, so Serre duality
gives the vanishing, and Riemann--Roch gives the dimension. The
same reasoning gives $h^0(L_i)=3$ and $h^1(L_i)=0$. Consequently
evaluation onto $(\mathcal E_i)_{q_i}$ is surjective, with a
one-dimensional kernel.

Put $V_i=H^0(E_i,L_i)$, and denote its dual evaluation image by
$U_i(q_i)\subset V_i^*$. It is a two-dimensional subspace. If
$q_i\ne q'_i$, the two evaluation kernels are different: otherwise
their common nonzero section would vanish on the degree-four
divisor $2p_i+2p'_i$, which a section of a degree-three line bundle
cannot do. Thus the two distinct planes in $V_i^*$ satisfy
\[
\dim\bigl(U_i(q_i)\cap U_i(q'_i)\bigr)=1.
\tag{DE-FI2}
\]

The product description of relative Frobenius and finite flat
base change give
$\mathcal E=\mathcal E_1\boxtimes\mathcal E_2$, while the
global section space is $V_1\otimes V_2$. The dual image of
evaluation at $q$ is therefore
$U_1(q_1)\otimes U_2(q_2)$. For subspaces $A,A'$ of one vector
space and $B,B'$ of another, adapted bases give the identity
\[
(A\otimes B)\cap(A'\otimes B')
                      =(A\cap A')\otimes(B\cap B').
\tag{DE-FI3}
\]
Both dual fibre images have dimension four. If both coordinates
are different, (DE-FI2)--(DE-FI3) show that their intersection has
dimension one, so their sum has dimension $4+4-1=7$.
If just one coordinate agrees, the intersection has dimension
$2\cdot1=2$, and the sum has dimension six. The rank of the dual
map equals that of (DE-FI1), proving the assertions. Relative
Frobenius twists change the target labels and scalar conventions,
but every map above is linear over the same algebraically closed
base after finite flat base change.
\end{proof}

\begin{remark}[Scope of the interpolation obstruction]
\label{rem:DE-A4-product-Frobenius-interpolation-scope}
\status{Proved scope; no generic-rank conclusion asserted.}
The lemma applies, in particular, to supersingular products and
to ordinary--supersingular products. It disproves the proposed
universal sufficient condition that $F_*L$ interpolate at
$(h^0(L)-1)/4$ general points when this number is integral. If a
Cartier pairing on $F_*L$ is alternating and a functional is a
sum of two fibre evaluations, its induced alternating form on
$H^0(L)$ factors through the rank-seven map (DE-FI1). An
alternating form on a seven-dimensional space has rank at most
six, so this particular construction cannot give rank eight.
No assertion about the rank for a general functional on
$H^0(L^{(2)})$, interpolation using more points, or the remaining
positive-index mixed-surface rank problem follows. The proof is
uniform linear algebra and elliptic-curve Riemann--Roch; no finite
computation is used or required.
\end{remark}
